% =================================================
% Учительская версия рабочего листа. Урок 3
% =================================================

\worksheettitle
  {Составление чисел. Задачи на цифры}
  {Урок 3 из 3. Версия для учителя}

\begin{teacherbox}[Цели урока]
  К концу урока ученик составляет все числа из заданных цифр, различает
  случаи с повторениями и без повторений, учитывает запрет ведущего нуля,
  строит упорядоченный перебор, находит наибольшее и наименьшее число и
  решает простые задачи на зависимость между цифрами.
\end{teacherbox}

\begin{checkpointbox}[Входная диагностика: ответы]
  \begin{enumerate}[label=\textbf{\arabic*.},itemsep=0.35em]
    \item Цифра десятков --- \(0\).
    \item \(5\,008\,002\).
    \item \(741\).
    \item Нет: \(053=53\), это двузначное число.
  \end{enumerate}
\end{checkpointbox}

\begin{teacherbox}[Как читать диагностику]
  Первые два задания возвращают к урокам 1--2. Если ученик игнорирует ноль
  в разряде десятков или записывает \(5\,802\), он ещё не удерживает
  позиционную запись и классы. Задания 3--4 проверяют готовность к новому
  материалу: порядок цифр и ограничение на старший разряд.
\end{teacherbox}

\section{Порядок цифр изменяет число}

\orderchangesfigure

\[
  27=2\text{ десятка}+7\text{ единиц},
  \qquad
  72=7\text{ десятков}+2\text{ единицы}.
\]

\begin{conclusionbox}
  Числа \(27\) и \(72\) составлены из одинаковых цифр, но цифры занимают
  разные разряды. Поэтому значения чисел различны.
\end{conclusionbox}

\begin{teacherbox}[Методический акцент]
  Не ограничивайтесь фразой «цифры переставили». Просите назвать значение
  каждой цифры: \(2\) означает \(20\) в числе \(27\), но только \(2\) в
  числе \(72\). Так новое содержание связывается с позиционным принципом.
\end{teacherbox}

\section{Цифры могут повторяться}

\repeatstreefigure

Если цифра десятков равна \(3\), получаем \(33\) и \(37\). Если цифра
десятков равна \(7\), получаем \(73\) и \(77\). Всего четыре числа:
\[
  \boxed{33,\ 37,\ 73,\ 77}.
\]

Для цифр \(2\) и \(5\):
\[
  \boxed{22,\ 25,\ 52,\ 55}.
\]

\begin{teacherbox}[Что значит «все числа»]
  Полнота списка --- отдельный учебный результат. Ученик должен не только
  назвать четыре числа, но и объяснить порядок: сначала все числа с
  десятками \(3\), затем все числа с десятками \(7\). Случайный список не
  даёт уверенности, что ответ полный.
\end{teacherbox}

\begin{teacherbox}[Уточняйте формулировку]
  Фраза «в записи входят только цифры \(3\) и \(7\)» допускает повторения:
  \(33\) и \(77\) подходят. Фраза «используя каждую цифру один раз» запрещает
  повторение. Эти условия нельзя считать равносильными.
\end{teacherbox}

\section{Цифры не повторяются}

\norepeattreefigure

Первая цифра не может быть равна \(0\). При первой цифре \(5\) получаем
\(508\) и \(580\); при первой цифре \(8\) --- \(805\) и \(850\).

\[
  \boxed{508,\ 580,\ 805,\ 850}.
\]

\[
  053=53,
\]
поэтому это двузначное, а не трёхзначное число.

Для цифр \(0,4,9\):
\[
  \boxed{409,\ 490,\ 904,\ 940}.
\]

\begin{teacherbox}[Типичная ошибка]
  Часть учеников записывает все шесть перестановок, включая \(049\) и
  \(094\). Не просто зачёркивайте их: предложите прочитать записи как числа.
  Получатся \(49\) и \(94\), то есть числа другой разрядности.
\end{teacherbox}

\section{Наибольшее и наименьшее число}

\largestsmallestfigure

Правила:
\begin{itemize}
  \item для наибольшего числа цифры располагают по убыванию;
  \item для наименьшего числа без нуля --- по возрастанию;
  \item при наличии нуля сначала ставят наименьшую ненулевую цифру, затем
    ноль и остальные цифры по возрастанию.
\end{itemize}

\begin{center}
\renewcommand{\arraystretch}{1.42}
\begin{tabular}{c|c|c}
  \toprule
  Цифры & Наибольшее & Наименьшее \\
  \midrule
  \(1,3,6,9\) & \(9631\) & \(1369\) \\
  \(0,2,4,7\) & \(7420\) & \(2047\) \\
  \(0,1,5,8,9\) & \(98510\) & \(10589\) \\
  \bottomrule
\end{tabular}
\end{center}

\begin{teacherbox}[Почему правило работает]
  Сравнение многозначных чисел начинается со старшего разряда. Поэтому для
  максимума в старший разряд помещают наибольшую цифру, затем наибольшую из
  оставшихся. Для минимума рассуждение обратное, но ноль не может занимать
  старший разряд.
\end{teacherbox}

\section{Условие связывает цифры}

\systematicroutefigure

Если цифра десятков на \(3\) больше цифры единиц:

\begin{center}
\renewcommand{\arraystretch}{1.35}
\begin{tabular}{c|ccccccc}
  \toprule
  цифра единиц & 0 & 1 & 2 & 3 & 4 & 5 & 6 \\
  \midrule
  цифра десятков & 3 & 4 & 5 & 6 & 7 & 8 & 9 \\
  \bottomrule
\end{tabular}
\end{center}

\[
  \boxed{30,\ 41,\ 52,\ 63,\ 74,\ 85,\ 96}.
\]

Цифра единиц \(7\) не подходит, поскольку цифра десятков должна была бы
равняться \(10\).

Если цифра единиц в два раза больше цифры десятков:

\begin{center}
\renewcommand{\arraystretch}{1.35}
\begin{tabular}{c|cccc}
  \toprule
  цифра десятков & 1 & 2 & 3 & 4 \\
  \midrule
  цифра единиц & 2 & 4 & 6 & 8 \\
  \bottomrule
\end{tabular}
\end{center}

\[
  \boxed{12,\ 24,\ 36,\ 48}.
\]

\begin{teacherbox}[Язык условия]
  В формулировке «цифра десятков на \(3\) больше цифры единиц» сначала
  зафиксируйте, какая цифра больше. Полезно проговорить равенство словами:
  «десятки = единицы + 3». Это предупреждает обращение зависимости.
\end{teacherbox}

\section{Считаем подходящие числа}

При сумме цифр \(7\):

\begin{center}
\renewcommand{\arraystretch}{1.35}
\begin{tabular}{c|ccccccc}
  \toprule
  десятки & 1 & 2 & 3 & 4 & 5 & 6 & 7 \\
  \midrule
  единицы & 6 & 5 & 4 & 3 & 2 & 1 & 0 \\
  \bottomrule
\end{tabular}
\end{center}

\[
  \boxed{16,\ 25,\ 34,\ 43,\ 52,\ 61,\ 70}.
\]
Всего \(7\) чисел.

\placevaluecountfigure

\begin{teacherbox}[Подсчёт и перечисление]
  На первом уроке по перебору список лучше сохранять даже тогда, когда
  требуется только количество. Сильным ученикам можно предложить объяснить
  ответ без полного списка: цифра десятков принимает семь значений, а цифра
  единиц после этого определяется однозначно.
\end{teacherbox}

\section{Повторяем разрядный состав}

\[
  5\text{ десятков тысяч}+78\text{ сотен}
  =5\cdot10\,000+78\cdot100
  =50\,000+7\,800
  =\boxed{57\,800}.
\]

\begin{teacherbox}[Диагностический смысл]
  Ошибка \(578\) или механическое приписывание \(5\) и \(78\) показывает,
  что ученик не переводит единицы счёта в их значения. Верните его к записи
  «один десяток тысяч = \(10\,000\)», «одна сотня = \(100\)».
\end{teacherbox}

\section{Самостоятельная проверка}

\begin{fillbox}[Ответы]
  \begin{enumerate}
    \item \(44,49,94,99\).
    \item \(207,270,702,720\).
    \item \(8630\) и \(3068\).
    \item \(20,31,42,53,64,75,86,97\).
  \end{enumerate}
\end{fillbox}

\begin{checkpointbox}[Выходной билет: ответы]
  \begin{enumerate}
    \item Начальный ноль не записывает сотни: \(041=41\).
    \item \(2059\).
    \item \(62\).
  \end{enumerate}
\end{checkpointbox}

\begin{teacherbox}[Итоговая оценка]
  Ученик достиг базовой цели, если он строит перебор по старшему разряду,
  отдельно проговаривает возможность повторения, исключает ведущий ноль и
  может обосновать полноту списка. Скорость перебора вторична по сравнению
  с его упорядоченностью.
\end{teacherbox}
